% Part: incompleteness
% Chapter: theories-computability
% Section: first-incompleteness

\documentclass[../../../include/open-logic-section]{subfiles}

\begin{document}

\olfileid{inc}{tcp}{inc}

\olsection{$\Th{Q}$ને પૂર્ણ, સુસંગત અને !!^{axiomatizable}
  વિસ્તરણો નથી}

\begin{thm}
\ollabel{thm:first-incompleteness}
$\Th{Q}$નું પૂર્ણ, સુસંગત અને
!!{axiomatizable} વિસ્તરણ
હોતું નથી.
\end{thm}

\begin{proof}
$\Th{Q}$નું કોઈ સુસંગત,
નિર્ણેય વિસ્તરણ નથી
તે આપણે પહેલેથી જાણીએ છીએ.
પરંતુ $\Th{T}$ પૂર્ણ અને
!!{axiomatized} હોય,
તો તે નિર્ણેય છે.
\end{proof}

\begin{explain}
આ પ્રમેય પ્રથમ અપૂર્ણતા
પ્રમેયની ગોડેલે ૧૯૩૧માં
આપેલી મૂળ રજૂઆતથી
ખૂબ દૂર નથી.
વધુ આધુનિક પરિભાષા સિવાય,
મુખ્ય તફાવત આ છે:
ગોડેલે ``સુસંગત''ને બદલે
``$\omega$-સુસંગત''
વાપર્યું હતું; અને સંગણનીયતાનો
ઔપચારિક ખ્યાલ ત્યારે હજી
ઉપલબ્ધ ન હોવાથી તેઓ
સંપૂર્ણ સામાન્ય અર્થમાં
``!!{axiomatizable}'' કહી
શકતા નહોતા. (પછીના
દાયકામાં, ગોડેલ સહિતના
સંશોધકોએ સંગણનીયતાના
ઔપચારિક નમૂનાઓ વિકસાવ્યા.
તેનું મોટું કારણ ગોડેલની
અપૂર્ણતાને આધીન થતા
સિદ્ધાંતોના પ્રકારો ઓળખવાનું
પણ હતું.)

પ્રમેય કહે છે કે પૂર્ણતા,
સુસંગતતા અને
!!{axiomatizability}
ત્રણેય એકસાથે મળી શકતાં
નથી. છતાં આમાંથી કોઈ
એક છોડીએ તો બાકીનાં બે
મળી શકે છે: $\Th{Q}$
સુસંગત અને સંગણનીય રીતે
સ્વયંસિદ્ધિત છે, પરંતુ
પૂર્ણ નથી; અસુસંગત
સિદ્ધાંત પૂર્ણ છે અને
સંગણનીય રીતે સ્વયંસિદ્ધિત
છે (દાખલા તરીકે
$\{ \eq/[0][0] \}$ વડે),
પરંતુ સુસંગત નથી; અને
અંકગણિતનાં સાચાં વાક્યોનો
ગણ પૂર્ણ અને સુસંગત છે,
પરંતુ સંગણનીય રીતે
સ્વયંસિદ્ધિત નથી.
\end{explain}
\end{document}
